\input amstex
\documentstyle{amsppt}
\magnification 1200
\pageno 10
\define\eps{\epsilon}
\define\okk{{\Cal O}_K}
\define\off{\overline{F}}
\define\ovf{\overline{F}}
\define\unb{\underline{B}}
\define\unt{\underline{T}}
\define\unx{\underline{x}}
\define\qqq{\Bbb Q}
\define\ccc{\Bbb C}
\define\zzz{\Bbb Z}
\define\rrr{\Bbb R}
\define\aaa{\Bbb A}
\define\nnn{\Bbb N}
\define\gga{\Bbb G}
\define\oee{{\Cal O}_E}
\define\VV{\Cal V}
\define\HH{\Cal H}
\hsize 6.5 true in
\vsize 8.7 true in
\vskip .3in
\subheading{\S 1 Some Background}
\vskip .1in
Let $G$ denote a connected semi-simple
quasi-split algebraic group defined over 
$F$ and let $T_G$ denote a maximal torus of $G$ defined
over $F$. 
In fact, we will assume that $T_G$ is an elliptic Cartan
but we will not make use of this immediately.
We shall work over a separable closure $\overline{F}$
most of this section.
A {\bf root} is a character
$\rho\in X^*(T)$ of a torus $T$ such that there exists an isomorphism
$\theta_\rho :\gga_a\rightarrow U_\rho$, $U_\rho\subset G$
a closed subgroup, for which $t^{-1}\theta_\rho(\xi)t=
\theta_\rho(\rho(t)\xi)$, for $t\in T$ and $\xi\in\gga_a$.
Here $\gga_a$ denotes the 1-dimensional additive affine group.
Let $B_G$ denote a Borel containing $T_G$,
$\Sigma_G$ a set of roots of $T_G$, $\Sigma_G^+$ the set of
positive roots with respect to $B_G$,
and $\Delta$ the set of simple roots in $\Sigma_G^+$.
When $G=G_2$, we usually write
$\Delta=\{\alpha,\beta\}$ and
$$
\Sigma_G^+=\{\alpha,\beta,\gamma=\alpha+\beta,
\delta=2\alpha+\beta,\epsilon=3\alpha+\beta,
\zeta=3\alpha+2\beta\}.
$$
\par
Fix a Borel $\underline{B}$ of $G$ defined over ${F}$, a Cartan 
$\underline{T}$ defined over $F$ contained in $\underline{B}$,
and a positive Weyl chamber $W_0$ of $T_G$ with respect to 
$\Sigma_G^+$.

Let $\Omega$ denote the absolute Weyl group of $T_G$.
We may occasionally identify the roots of $G$ with respect to $T_G$
with those of $G$ with respect to $\underline{T}$ by means of
a fixed isomorphism $T_G\rightarrow \underline{T}$, and, thus, identify
the absolute Weyl group of $\underline{T}$ with that of $T_G$.
Representing $\omega\in \Omega$ by
$w\in N_{G(\overline{F})}({T}_G(\overline{F}))$,
define
$$
{B}_G(W):=w^{-1}{B}_Gw={B}_G^w,
$$
for $W=\omega^{-1}W_0$. 
The Borels
$\{{B}_G(W))\ |\ W\in {\Cal W}\}$ are precisely the Borel
subgroups of $G$ which contain ${T}_G$, where ${\Cal W}$ denotes
the set of Weyl chambers of $T_G$.
\par
Recall that a {\bf regular star} is a $|\Omega|$-tuple of Borels
$(B(W))_{W\in {\Cal W}}$ such that, for all $W\in {\Cal W}$,
$B(W)={B}_G(W)^h=h^{-1}{B}_G(W)h$, for some
$h\in G(\overline{F})$ (independent of $W$).
The variety of regular stars is denoted by $S_0$.
There is a natural bijection between $T_G(\ovf)\setminus G(\ovf)$
and $S_0(\ovf)$ given by $g\longmapsto (B(W)^g)_{W\in {\Cal W}}$.
\par
Let
$$
X_0:=\{(g,(B(W))_{W\in {\Cal W}})\in G\times S_0\ |\  
g\in\cap_{W\in {\Cal W}}B(W)\text{ is regular}\}, 
\tag 1.1
$$
so that we have an $\overline{F}$-isomorphism
$$
T_G'\times T_G\setminus G
\rightarrow X_0
\tag 1.2
$$
$$
(t,gT_G)\longmapsto (g^{-1}tg,{B}_G(W)^g).
$$
\par
\proclaim{Definition} Let
$\phi:X_0\rightarrow T_G$ be the map which sends 
$(g,s)\in {\Cal C}(T_G)\times S_0$ to the point in $T_G$ which meets the
conjugacy class of $g$, where ${\Cal C}(T_G)$ denotes the 
union of the conjugates of $T_G$. 
Let $\pi :X_0\rightarrow G$ denote the projection map onto
the ``$G$-coordinate''.
\endproclaim
The map $\pi$ restricts to an isomorphism onto the
conjugacy class of $\gamma$,
$$
\pi:\phi^{-1}(\gamma)\rightarrow \{g\gamma g^{-1}\ |\ g\in G/T_G\}.
$$
\vskip .3in
\subheading{Construction of $Y_1$}
\vskip .1in
Let $C\subset T_G$ be a one-parameter subgroup of $G$
containing $\{1\}$ but no other singular point.
Let $C':=C-\{1\}$
and let $Y_0:=\phi^{-1}(C')\subset X_0$. The isomorphism
(1.2) restricts to an $\overline{F}$-isomorphism
$$
C'\times T_G\setminus G\rightarrow Y_0.
\tag 1.4
$$
It can be shown that $Y_0$ is a regular algebraic subvariety of $G\times S$
\cite{L, p. 480} but is not complete. 
Let $Y$ denote the closure of $Y_0$ in $G\times S$. $Y$ is
singular in general.
\par
Let ${\Cal B}\cong B_G\setminus G$ denote the Borel variety.
It will be helpful
later to record here local coordinates for ${\Cal B}$.
First we recall the following 
well-known result \cite{Bor, ch IV, Corollary 14.13}:
\proclaim{Lemma 1.5} If $B,B',B''\in {\Cal B}$ and $B', B''$ are opposite
to $B$ then they are conjugate by an element of $B$.
\endproclaim
 As a consequence of the Levi decomposition, 
it follows that, for $B, B',B''$ as in the lemma,
$B'$ and $B''$ are conjugate by an element of $N$, the unipotent radical of 
$B$. Let ${\Cal B}(B)\subset {\Cal B}$
denote the open subvariety of all Borels
opposite to $B$.
\proclaim{Corollary 1.6} For each $B\in {\Cal B}$, ${\Cal B}(B)$
is isomorphic
as an affine variety to $N$ and the coordinates on $N$ (i.e., the matrix
coefficients) serve as coordinates on ${\Cal B}(B)$.
\endproclaim
A point
$(B(W))_{W\in{\Cal W}}\in {\Cal B}^{|\Omega|}$ is called a {\bf star}
if for every $\omega\in\Omega$ and every simple root $\alpha$
of $T_G$ in $B_G$ there is an $h\in G(\overline{F})$
such that
$$
B(\omega^{-1}W_0)=h^{-1}B_Gh
$$
and
$$
B((s_\alpha\omega)^{-1}W_0)\subset h^{-1}P_{G,\alpha} h
$$
where ${P}_{G,\alpha}$ is the smallest parabolic subgroup 
of $G$ containing $B_G$ and $B_G^{s_\alpha}$.
Here $s_\alpha$ denotes the reflection associated to
the simple root $\alpha$ of $T_G$ in $B_G$. The variety of stars,
denoted $S$, is closed in ${\Cal B}^{|\Omega |}$, projective,
but not in general regular as a variety.
It is relatively easy to see that a regular star is a star.
In fact, the variety of stars is the Zariski closure of $S_0$
in ${\Cal B}^{|\Omega|}$. (One way to see this is by means of
a more explicit realization of $S_0$ as a subset of $S$ using
coordinate patches (Lemma 1.2 below). That $S_0$ is 
open in its Zariski closure is a
well-known result in group theory, e.g.,
\cite{Sp3, ch.4, Lemma 4.3.1}, \cite{Bor, II.6, Prop. 6.7}.)
\par
We now define some open patches on $S$ and regular
coordinates on these patches.
Fix a Borel $B_+$ of $G$ and another Borel $B_-$ which is 
opposite to $B_+$; let $T_+:=B_+\cap B_-$. Let
$$
S(B_+):=\{(B(W))_{W\in{\Cal W}}\in S\ |\ \forall W\in{\Cal W},\ 
B(W) \text{ is opposite to }B_+\}
$$
and let
$$
S(B_+,B_-):=\{(B(W))_{W\in{\Cal W}}\in S(B_+)\ 
|\ B(W_0)=B_-\}.
$$
Let $N_+$ denote the unipotent radical of $B_+$. The map
$$
((B(W))_{W\in{\Cal W}},n)\longmapsto (B(W)^n)_{W\in{\Cal W}}
$$
defines an isomorphism
$$
S(B_+,B_-)\times N_+ \rightarrow S(B_+).
\tag 1.7
$$
This is because for $B_-$ and $B_+$ as above, the map
$n\longmapsto B_-^n$, $n\in N_+$, defines an isomorphism from
$N_+$ to the open subvariety ${\Cal B}(B_+)\subset {\Cal B}$
of all Borels opposite to $B_+$.
\par
Following \cite{L}, we now describe local coordinates for 
$S(B_+,B_-)$.
\par
We identify the absolute Weyl group of $T_G$ with the absolute
Weyl group of
$T_+$.
We identify the set of Weyl chambers of $T_G$ with those of $T_+$.
If $(W,\rho)$ denotes a Weyl chamber and a 
root defining a wall
of it (which we refer to below as a {\bf wall-chamber pair}
or often simply as a {\bf pair}), 
we may find a $\omega\in\Omega$ for which
$W=\omega^{-1}W_0$ and $\rho=\omega \alpha$, for some 
simple root $\alpha$ of $T_+$ in $B_-$;
we sometimes write $\alpha(W,\rho)$ for
the $\alpha$ obtained in this way.
\par
Given a star $(B(W))_{W\in{\Cal W}}\in S(B_+,B_-)$,
to each pair $(W,\rho)$ one defines
the value at $(B(W))_{W\in{\Cal W}}$ of the coordinate
function $z(W,\rho)$ as follows. Let $ker(\rho)
\subset X^*(T_G)\otimes {\rrr}$
denote the wall separating two adjacent Weyl chambers
$W$, $W'$ and let $h\in N_+(\ovf)$ denote the unique element
satisfying 
$$
B(W)=h^{-1}B_-h
$$
and
$$
B(W')\subset h^{-1}P_{-,\alpha} h
$$
where $P_{-,\alpha}$ is the smallest parabolic subgroup 
of $G$ containing $B_-$ and $B_-^{s_\alpha}$, $\alpha 
=\alpha (W,\rho)$.
Recall that
$$
\gga_a\cong U_{-\alpha}\cong P_{-,\alpha}/B_-,
$$
via $z\longmapsto u=exp(X_{-\alpha})\longmapsto B_-^u$, which we may
regard as an element of $P_{-,\alpha}/B_-\hookrightarrow G/B_-\cong {\Cal B}$.
Let ${\Bbb A}$ denote the affine line.
There is a coordinate function $z(W,\rho)$,
$$
z(W,\rho):S(B_+,B_-)\rightarrow \aaa,
$$
such that
$$
hB(W')h^{-1}=\theta_{-\alpha} (z(W,\rho))B_-\theta_{-\alpha} (-z(W,\rho)),
\ \ \alpha=\alpha(W,\rho),
\tag 1.8
$$
where $X_\alpha$, $X_{-\alpha}$ are root vectors satisfying
$$
[X_\alpha, X_{-\alpha}]=H_\alpha,\ \  
\beta(H_{\alpha})=2{(\beta,\alpha)\over (\alpha,\alpha)},\ \ 
\theta_{-\alpha}(z)=\exp(zX_{-\alpha})\in U_{-\alpha}.
$$
The local coordinate $z(W_0,\alpha)$ on $S(B_+,B_-)$ is
a local coordinate on a line of type $\alpha$ in the sense
of \cite{St}. This observation will be useful when dealing with the
comparison of the subregular germ.

We have

(a) $z(W,\rho)=0$ if and only if $B(W)=B(W')$ and 

(b) the collection of the $z(W,\rho)$, for all pairs $(W,\rho)$,
uniquely determines the star $(B(W))_{W\in{\Cal W}}\in S(B_+,B_-)$,

(c) the coordinate functions $z(W,\rho)$ are regular and generate
the coordinate ring of $S(B_+,B_-)$.

The coordinates $z(W,\rho)$ are not algebraically independent.
The coordinate relations they satisfy can be most naturally
described using the notion of a gallery in the spherical building 
associated to the Weyl group $\Omega$ of $T_G$. We may identify the
spherical building of $T_G$ with that of $T_+$.
Recall that a {\bf gallery} is a sequence $W_0,W_1,\cdots, W_r$
of Weyl chambers such that $W_i$, $W_{i+1}$ are distinct but share a wall
$ker(\rho_{i+1})$. We always take $W_0$ to be the positive Weyl chamber of
$T_+$ with respect to $B_+$ and define
$\omega_i\in\Omega$ to be that element for which
$W_i=\omega_i^{-1}W_0$. Let $\alpha_{i+1}:=\omega_i\rho_{i+1}$,
so $\alpha_{i+1}=\alpha(W_i,\rho_{i+1})$. If $z_i:=z(W_{i-1},\rho_i)$
for $1\leq i\leq r$,
then
$$
B(W_{i-1})=h^{-1}B_-h,\ \ \ 
B(W_i)=h^{-1}\theta_{-\alpha_i}(-z_i)
B_-\theta_{-\alpha_i}(z_i)h,
\tag 1.9
$$
for some $h\in N_+(\overline{F})$ possibly depending on $i$. We
define $n,n_{W_j}$, for $1\leq j\leq r$, by
$B(W_j)=B_-^{n_{W_j}n}$, where $n,n_{W_j}\in N_+(\overline{F})$
and
$$
n_{W_j}=\exp (z_jX_{-\alpha_j})
\exp (z_{j-1}X_{-\alpha_{j-1}})\cdots \exp (z_1X_{-\alpha_1}).
\tag 1.10
$$
In particular, if $W_r=W_0$ (so that the gallery is closed),
then
$$
\exp (z_rX_{-\alpha_r})
\exp (z_{r-1}X_{-\alpha_{r-1}})\cdots \exp (z_1X_{-\alpha_1})=1.
\tag 1.11
$$
In particular, if $W$, $W'$ are adjacent and share a wall
$ker(\rho)$ then $z(W,\rho)=-z(W',\rho)$.
\par
We now state an observation of Langlands \cite{L, pp.483-484}
\proclaim{Lemma 1.12} TFAE:
\newline
(a) A star $(B(W))_{W\in{\Cal W}}\in S(B_+,B_-)$ is regular.
\newline
(b) $B(W)\not= B(W')$ for all adjacent $W$, $W'$.
\newline
(c) $z(W,\rho)\not= 0$ for all pairs $(W,\rho)$.
 \endproclaim

\par
We define coordinate patches on a resolution $S_1$ of $S$ (\cite{H1}).
\par
Recall that all the short (resp. long) roots lie in the same orbit
under the Weyl group. With this fact and the above lemma in mind, define
$$
S'(B_+,B_-):=\{(B(W))_{W\in{\Cal W}}\in S(B_+,B_-)
\ |\ \forall \alpha\in\Delta ,\exists \omega\in\Omega
\text{ s.t. }z(\omega^{-1}W_0,\omega\alpha)\not= 0\}.
$$
Suppose $(B(W))_{W\in{\Cal W}}\in S(B_+,B_-)$ is not in
$S'(B_+,B_-)$. Then $z(\omega^{-1}W_0,\omega\alpha)=0$
for all $\alpha\in\Delta$ and $\omega\in\Omega$. By the lemma,
this implies $B(W)=B(W')$ for all adjacent Weyl chambers. This is 
the diagonal in ${\Cal B}^{|\Omega|}$ and the diagonal contains no
regular stars.
\par
Let $\aaa^\Delta$ denote the $|\Delta|$-fold Cartesian product
of affine spaces and let $T_+$ act on $\aaa^\Delta$ by
$$
t:(z(\alpha))_{\alpha\in\Delta}\longmapsto
(\alpha(t)^{-1}z(\alpha))_{\alpha\in\Delta},
$$
and on $S'(B_+,B_-)$ by
$$
t:(z(W,\rho))_{(W,\rho)\text{ a pair}}\longmapsto 
(\alpha(t)z(W,\rho))_{(W,\rho)\text{ a pair}},\ \ \ \alpha=\alpha(W,\rho).
$$
Let
$$
S_1(B_+,B_-):=S'(B_+,B_-)\times_{T_+}\aaa^\Delta:=
(S'(B_+,B_-)\times\aaa^\Delta)/T_+.
\tag 1.13
$$
Observe that there is a ``blow-down'' map
$$
S_1(B_+,B_-)\rightarrow S(B_+,B_-),
\tag 1.14
$$
given by $(s',z)T_+\longmapsto s$, where 
$s'=(z_1(W,\rho))_{(W,\rho)\text{ a pair}}$, 
$z=(z(\alpha))_{\alpha\in\Delta}$ (note $z(\alpha)$ is only defined up to
multiplication by $\alpha(t)^{-1}$), and
$$
s=(z(W,\rho))_{(W,\rho)\text{ a pair}}
=(z(\alpha(W,\rho))z_1(W,\rho))_{(W,\rho)\text{ a pair}}.
$$
This is clearly well-defined (i.e., independent of the chosen
representative
$z=(z(\alpha))_{\alpha\in\Delta}$ of the $T_+$-orbit in $\aaa^\Delta$).
The subscript 1 on $z_1(W,\rho)$ has been added to distinguish it
from its image 
$z(W,\rho):=z(\alpha(W,\rho))z_1(W,\rho)$ in $S(B_+,B_-)$.
Later we shall fix a canonical choice of a generic representative
$(z(\alpha))_{\alpha\in\Delta}$ of the $T_+$-orbit in $\aaa^\Delta$:
$z(\alpha)$ is chosen to be that multiple of $\alpha(t)^{-1}$
such that $z(\alpha)=\lambda /x(\alpha)$, where $x(\alpha)$ is as in
(1.17) below (equivalently: $z(\alpha)$ is chosen to be that multiple
which makes $w(\alpha)=1$ for $\alpha >0$ simple \cite{H1, ch II, \S 2}).
\par
Let $S'(B_+)\subset S(B_+)$ be the image of
$S'(B_+,B_-)\times N_+$ in $S(B_+)$ under the map
$$
((B(W))_{W\in{\Cal W}},n)\longmapsto (B(W)^n)_{W\in{\Cal W}},
$$
and let $S_1(B_+)\subset S'(B_+)\times_{T_+}\aaa^\Delta$ 
be the image of $S_1(B_+,B_-)\times N_+$ under
$$
((B(W))_{W\in{\Cal W}},(z(\alpha))_{\alpha\in\Delta},n)
\longmapsto ((B(W)^n)_{W\in{\Cal W}},(z(\alpha))_{\alpha\in\Delta}).
$$
Therefore,
$$
S'(B_+)\cong S'(B_+,B_-)\times N_+,
\tag 1.15a
$$
and 
$$
S_1(B_+)\cong S_1(B_+,B_-)\times N_+.
\tag 1.15b
$$
The apparent dependence of $S_1(B_+)$ on $B_-$ is not serious:
by \cite{L, Lemma 3.6}, if we replace $B_-$ by another
Borel $B'_-$ opposite
to $B_+$ then the corresponding variety is canonically isomorphic
to $S_1(B_+)$.
Let $B'_+$ be another Borel and let $B_-$ be a Borel opposite to both
$B_+$ and $B'_+$. There is a birational morphism from
$S_1(B_+)$ to $S_1(B'_+)$ as defined in \cite{L}. Using this birational map
to paste $S_1(B_+)$ and $S_1(B'_+)$ together, we obtain
a variety $S_1$.
\par
Observe that the ``obvious'' map $S_0(B_+)\rightarrow S_1(B_+)$
extends to a morphism $S_0\rightarrow S_1$. This map yields a map
$X_0\rightarrow G\times S_1$.
\proclaim{Definition 1.16}
Let $X$ be the Zariski closure of $X_0$ in $G\times S$.
Let $X_1$ be the closure of $X_0$ in $G\times S_1$ and let $Y_1$
be the closure of $Y_0$ in $X_1$. Let $X_1(B_+)$ denote the
intersection of $X_1$ with $G\times S_1(B_+)$ and similarly for
$Y_1(B_+)$. Similarly, we define $X''$, $Y''$, $X''(B_+)$,
$Y''(B_+)$, replacing $G\times S_1(B_+)$
by $G\times S''$ or $G\times S''(B_+)$ in the definition of
$X_1$ and $Y_1$, where $S''$ is the image of $S_0$ in $Y_1$.
\endproclaim
In particular, the first coordinate of a point
$p\in Y''(B_+,B_-)\subset G\times S''$ is an element
$b$ of $\cap_W B_-^{n_W} \subset B_-$, which we may write as
an ordered product
$$
b=t_0\prod_{\rho\in\Sigma_G^+}\exp (x(\rho)X_\rho),
\tag 1.17
$$
where $\Sigma_G^+$ denotes the set of roots of $T_+$
which are positive with respect to the Borel $B_-$, on which
we fix an ordering once and for all.
\vskip .3in
\subheading{Coordinates on $Y''$}
\vskip .1in
Consider a point $x\in X_1\subset G\times S_1$.
On the patch of $X_1$ meeting
$G\times S_1(B_+,B_-)\times N_+\cong
G\times S_1(B_+)$, we write $x$ as 
$(b,(B(W))_{W\in{\Cal W}},(z(\alpha))_{\alpha\in\Delta},n)$,
where $b\in\cap_{W\in{\Cal W}}B(W)\subset B_-$ and
$B(W_0)=B_-$ (the condition that $b\in\cap_{W\in{\Cal W}}B(W)$
is inherited from the definition of $X_0$).
\par
Other coordinates on $X''(B_+)$ may be defined in terms of a point
$(t^g,(B(W)^g)_{W\in{\Cal W}})$ in 
\newline
$G\times S(B_+,B_-)$, where $t\in T_G(\ovf )$ is regular. We may write
$$
g=g_0t_+n_-n_+,
$$
where $T_G^{g_0}=T_+$, $t_+\in T_+(\ovf)$, $n_-\in N_-(\ovf)$,
and $n_+\in N_+(\ovf)$. Let $t^{g_0}=t'_+$, for some
$t'_+\in T_+(\ovf)$, and write $t^g=(t'_+n'_-)^{n_+}$,
for some $n'_-\in N_-(\ovf)$. Therefore,
$$
(t^g,(B(W)^g)_{W\in{\Cal W}})=(t'_+n'_-,B_-^{n_W})^{n_+},
$$
where $n_W\in N_+(\ovf)$. We write
$$
n'_-=\prod_{\rho\in\Sigma_G^+}\exp (x(\rho)X_\rho),
$$
$$
n_+=\prod_{\rho\in\Sigma_G^+}\exp (v(\rho)X_{-\rho}),
$$
where, by (1.10),
$$
n_W=\exp (z(W_r,\rho_r)X_{-\alpha_r})
\exp (z(W_{r-1},\rho_{r-1})X_{-\alpha_{r-1}})\cdots 
\exp (z(W_1,\rho_1)X_{-\alpha_1}),
\tag 1.18
$$
where $W=\omega^{-1}W_0$, $\omega=s_{\alpha_r}\cdots s_{\alpha_1}$,
and $W_0$, ..., $W_r$ being chambers in a path from the positive Weyl chamber
$W_0$
to $W_r=W$. Recall we defined $z_1(W,\rho)$, for any pair $(W,\rho)$, by
$$
z(W,\rho)=z(\alpha(W,\rho))z_1(W,\rho).
\tag 1.19
$$
\par
Define the coordinate function $w(\rho)$, $\rho >0$, by
$w(\alpha)=1$, if $\alpha >0$ is simple, and otherwise by
$$
(t'_+)^{-1}(t'_+)^{n_-}=
\prod_{\rho\in\Sigma_G^+}
\exp ({w(\rho)\lambda\over
\prod_{\alpha\in\Delta}z(\alpha)^{m(\alpha)}}X_\rho),
\tag 1.20
$$
where $\rho=\sum_{\alpha\in\Delta}m(\alpha)\alpha$. This describes a point
in the image
$$
C'\times G/T_G\hookrightarrow Y_1(B_+,B_-)
\subset X_1(B_+,B_-)\subset G\times S_1(B_+,B_-),
$$
where $C'=C\cap T_G'=C-\{1\}$. 
\vskip .3in
{\bf Coordinates on $Y_1$}
\vskip .1in
We begin with the following
\proclaim{Lemma 1.21}(\cite{H1,ch II, Lemma 3.3}) If $t\in T_+$ and
$$
n=\prod_{\rho >0}\theta_\rho(n_\rho)\in N_+
\ \ \ \ \ (\text{ordered product}),
$$
then
$$
t^n=t\prod_{\rho >0}\theta_\rho(y(\rho)),
$$
where
$$
y(\rho)=\sum_{\rho_1,...,\rho_m} c_{\rho_1,...,\rho_m}(t)
n_{\rho_1}...n_{\rho_m},
$$
the sum ranging over all roots $\rho_i>0$ such that
$\rho=\rho_1+...+\rho_m$ and
$$
c_{\rho_1,...,\rho_m}(t)=
\sum_{\sigma_1,\sigma_2}(1-\sigma_1(t)^{-1})\sigma_2(t),
$$
where the sum over $\sigma_1,\sigma_2$ such that $\sigma_1$ is a root
and $\sigma_2$ is a sum of roots. Moreover, $c_\rho(t)=(1-\rho(t)^{-1})$.
\endproclaim
We want to apply this lemma to obtain coordinates of a point
$$
(t^g,(B(W))_{W\in{\Cal W}},(z(\alpha))_{\alpha\in\Delta})^n\in
\chi^{-1}(C)\times S_1(B_+,B_-)\times N_+.
$$
For the case $G=G_2$, we
give coordinate relations between the coordinates in (1.20) and
the $n_\rho$'s and the $z(W,\alpha)$'s (see \cite{H1,chap II}):
$$
x(\alpha)=(1-\alpha^{-1})n_\alpha = (1-\alpha^{-1})/z(W_0,\alpha),
\tag 1.22
$$
$$
x(\beta)=(1-\beta^{-1})n_\beta = (1-\beta^{-1})/z(W_0,\beta),
$$
$$
x(\gamma)=(1-\gamma^{-1})n_\gamma +
c_1\cdot {(1-\alpha^{-1})
\over{z(W_0,\alpha)z(W_0,\beta)z(\sigma_\beta^{-1}W_0,\alpha)}}
+\text{ higher order terms in }\lambda,
$$
$$
x(\delta)=(1-\delta^{-1})n_\delta +
c_2\cdot {(1-\alpha^{-1})
\over{z(W_0,\alpha)z(W_0,\beta)z(\sigma_\beta^{-1}W_0,\alpha)}}
+\text{ higher order terms in }\lambda,
$$
$$
x(\epsilon)=(1-\epsilon^{-1})n_\epsilon +
c_3\cdot {(1-\alpha^{-1})n_\delta
\over{z(W_0,\alpha)}}
+\text{ higher order terms in }\lambda,
$$
$$
\split
x(\zeta)
&=(1-\zeta^{-1})n_\zeta +
c_4\cdot {(1-\alpha^{-1})n_\gamma^2
\over{z(W_0,\alpha)}}+
c_5\cdot {(1-\alpha^{-1})n_\delta
\over{z(W_0,\alpha)z(W_0,\beta)}}+\\
&+c_6\cdot {(1-\beta^{-1})n_\epsilon
\over{z(W_0,\beta)}}+
c_7\cdot {(1-\gamma^{-1})n_\delta
\over{z(W_0,\beta)z(\sigma_\beta^{-1}W_0,\alpha)}}
+\text{ higher order terms in }\lambda,
\endsplit
$$
$$
\split
w(\gamma)
&={x(\gamma)z(\beta)z(\alpha)\over \lambda}\\
&=(\gamma(X)n_\gamma + c_1{\alpha(X)\over{z(W_0,\alpha)z(W_0,\beta)}})
\alpha(X)^{-1}\beta(X)^{-1}z(W_0,\alpha)z(W_0,\beta)+\\
&\quad\quad +\text{ higher order terms in }\lambda,
\endsplit
$$
$$
\split
w(\delta)
&={x(\delta)z(\beta)z(\alpha)^2\over \lambda}\\
&=(\delta(X)n_\delta +
c_2{\alpha(X)\over{z(W_0,\alpha)z(W_0,\beta)z(\sigma_\beta^{-1}W_0,\alpha)}})
\alpha(X)^{-2}\beta(X)^{-1}z(W_0,\alpha)^2z(W_0,\beta)+\\
&\quad\quad +\text{ higher order terms in }\lambda,
\endsplit
$$
$$
\split
w(\epsilon)
&={x(\epsilon)z(\beta)z(\alpha)^3\over \lambda}\\
&=(\epsilon(X)n_\epsilon + c_3{\alpha(X)n_\delta\over{z(W_0,\alpha)}})
\alpha(X)^{-3}\beta(X)^{-1}z(W_0,\alpha)^3z(W_0,\beta)+\\
&\quad\quad +\text{ higher order terms in }\lambda,
\endsplit
$$
$$
\split
w(\zeta)
&={x(\zeta)z(\beta)^2z(\alpha)^3\over \lambda}\\
&=(\zeta(X)n_\zeta + c_4{\alpha(X)n_\gamma^2\over{z(W_0,\alpha)}}
+c_5{\alpha(X)n_\delta\over{z(W_0,\alpha)z(W_0,\beta)}}
+c_6{\beta(X)n_\epsilon\over{z(W_0,\beta)}}+\\
&\quad +{\gamma(X)n_\zeta\over{z(\sigma_\beta^{-1}W_0,\alpha)z(W_0,\beta)}})
\alpha(X)^{-3}\beta(X)^{-2}z(W_0,\alpha)^3z(W_0,\beta)^2+\\
&\quad\quad +\text{ higher order terms in }\lambda,
\endsplit
$$
Because of (1.22), on $Y_1(B_+,B_-)$, we may
use the algebraically independent coordinates
$$
z(W_0,\alpha),\ z(W_0,\beta),\ z(\sigma_\beta^{-1}W_0,\alpha),\
n_\gamma,\ n_\delta,\ n_\epsilon,\ n_\zeta,
$$
in place of the coordinates
$$
z(\alpha),\ z(\beta),\ x(\alpha),\ x(\beta),\ x(\gamma),\
x(\delta),\ x(\epsilon),\ x(\zeta),\
w(\gamma),\ w(\delta),\ w(\epsilon),\ w(\zeta),
$$
on $Y''(B_+,B_-)$ appearing in (1.17) and (1.20).

\enddocument
